\documentstyle[%


righttag%


,11pt%


,amssymb%


,russian%               remove in Eng.


,izvru%                 Set izven% in Eng.


]{amsart}





%       Math definitions





\newcommand{\re}{\operatorname{Re}}


\newcommand{\tts}[1]{}


\numberwithin{equation}{section}


\renewcommand{\baselinestretch}{0.96}





%\hoffset=-15mm%                  For Eng.


\setcounter{page}{23}


\god{2000}


\nomer{5~(456)} %                        Remove in Eng.


%\nomer{Vol.\,44, No.\,5}%               Set in Eng.


%\str{\pageref{begin}--\pageref{end}}%  Set in Eng.


\udk{519.216}





\author{\it Л.А.\,Золотухина}


% NB:   ^^ remove in Eng.


\title{Асимптотическое распределение числа


совпадений двумерной выборки при


натуральном совмещении}





\date{}


%\pagestyle{myheadings}%         Set in Eng.


\pagestyle{plain} %              Remove in Eng.


\begin{document}


%\thispagestyle{plain}%          Set in Eng.


\maketitle


\label{begin}





\section{Введение}





Классическая задача о совпадениях сформулирована Монмортом в 1708 г.


Пусть мы имеем, например, две колоды карт. Перетасуем колоды каждую в


отдельности, а затем разложим на столе одну под другой. Если сверху и


снизу оказалась одна и та же карта, то будем говорить о совпадении.


Известно, что среднее число совпадений при таком сопоставлении при любом


числе карт в колоде равно 1. Если же будем увеличивать число карт в


колоде, распределение числа совпадений будет стремиться к распределению


Пуассона с параметром 1.





В работах [1]--[3] было рассмотрено следующее обобщение этой задачи.


Пусть выборка $(u_i,t_i)$ ($i=1,\dotsc,n$) объема $n$ извлечена из двумерной


генеральной совокупности с плотностью распределения $f(x,y)>0$ при


$a<x<b$, $c<y<d.$ Возможно $a$, $c$ равны $-\infty$, a $b$, $d$ равны


$+\infty$.


Но компоненты каждого выборочного вектора $(u_i,t_i)$ наблюдаются отдельно,


так что имеем два набора наблюдений:





$\Vec t=(t_1,\dotsc ,t_n)$ --- наблюдения первой компоненты,





$\Vec u=(u_1,\dotsc ,u_n)$ --- наблюдения второй компоненты.





Затем элементы каждой из выборок располагаются в случайном порядке (любая


перестановка $(t_1, \dotsc t_n)$ и $(u_1, \dotsc u_n)$ является


равновероятной). Необходимо найти такой способ


соединения элементов выборок $\Vec t $ и $\Vec u $  в пары, чтобы наилучшим


образом репродуцировать


пары исходной выборки. B [1] было предложено несколько критериев


оптимальности, в частности, а)~максимизация вероятности репродукции пар и


б)~максимизация среднего числа совпадений.


Способ соединения в пары, удовлетворяющий критерию~б), был назван


``оптимальным'' в отличие от удовлетворяющего критерию~а) ``максимального


правдоподобия'' (ML).


Способ, при котором в пары объединяются соответствующие элементы


вариационных рядов, будем называть натуральным совмещением.





В совокупности в [1]--[3] были получены следующие результаты.


Пусть плотность распределения генеральной совокупности


$f(x,y)>0$, $x \in R$, $y\in R$, и обладает монотонным


отношением правдоподобия


$$


f(x_1,y_1)f(x_2,y_2) \ge f(x_1,y_2)f(x_2,y_1)\ \;


\text{для} \ \; x_1<x_2; \ \; y_1<y_2.


$$


Тогда а)~ML-правило есть натуральное совмещение;


б)~оптимальное правило предполагает объединение в пары наибольших и


наименьших элементов $\Vec t$ и $\Vec u$;


в)~при натуральном совмещении среднее число совпадений  не меньше 1, т.\,е.


натуральное совмещение всегда не хуже случайного объединения в пары.





Все эти выводы были сделаны авторами для любых объемов выборки $n$. Позже


в [4], [5] были получены интегральные


представления математического ожидания и дисперсии числа совпадений, а


также вычислена асимптотика при $n \to +\infty$ среднего числа совпадений


при натуральном совмещении.





Целью данной статьи является получение интегрального представления всех


моментов числа совпадений (п.2), а также доказательство того факта, что


асимптотическое распределение числа совпадений --- распределение Пуассона с


параметром


$\lambda={\displaystyle\int_{a}^{b}}\dfrac{f(x,y^\ast (x))}


{f_2(y^\ast (x))}dx<\infty$,


где $f_2(y)=\displaystyle\int_a^b{f(x,y)}dx>0$ --- маргинальная


плотность распределения второй компоненты,


$$


F_1(x)=\int_a^x\int_c^df(x,y)dx\,dy,\quad


F_2(x)=\int_a^b\int_c^yf(x,y)dx\,dy


$$


--- маргинальные функции распределения $u_i$ и $t_i$ соответственно, а


$y^*(x)=F_2^{-1}(F_{1}(x))$.





Если же $\lambda=\infty,$ то распределение числа совпадений асимптотически


нормально, точнее


$$


P\bigg(\frac{N-MN}{\sqrt{DN}} <x\bigg) \rightarrow \phi(x)=


\frac{1}{\sqrt{2\pi}}


\int_{-\infty}^x \exp \bigg(\frac{-u^2}{2}\bigg)du.


$$





Кроме того, п.\,5 содержит доказательство следующего утверждения: для


любого $n$ число совпадений распределено по рандомизированному биномиальному


закону.





\section{Интегральные представления моментов}





В дальнейшем для определенности положим $f(x,y)>0$, $x \in R$, $y \in R$.


Итак, пусть из двумерной генеральной совокупности с плотностью $f(x,y)$


извлечена выборка $((t_1, \dotsc, t_n),(u_1, \dotsc, u_n))$ отдельно по


компонентам. Пусть, кроме того, $(x_1,x_2, \dotsc, x_n)$ --- упорядоченные


$(t_1,t_2, \dotsc, t_n)$, а $(y_1,y_2, \dotsc, y_n)$ --- упорядоченные


$(u_1,u_2, \dotsc, u_n)$. Будем объединять в пары $x_i$ c $y_i$ --- в этом


состоит натуральное совмещение.





Случайные величины $\chi_i$ введем следующим образом:


$\chi_i=1$, если среди образовавшихся пар присутствует пара $(t_i,u_i)$, и


$\chi_i=0$ в противном случае.





Пусть $N$ --- число совпадений.


Очевидно, $N=\sum\limits_{i=1}^n\chi_i$ и


$\chi_i$ зависимы симметрично. Тогда


\begin{equation}


MN^k=M\bigg(\sum_{i=1}^n\chi_i\bigg)^k=


M\sum_{(k_1,\dotsc,k_n)}\frac{\chi_1^{k_1}\chi_2^{k_2}


\dotsb \chi_n^{k_n}}{k_1!k_2! \dotsb k_n!}k! =


\sum_{j=1}^k \frac{c_{jk} n!}{(n-j)!}


M(\prod_{m=1}^j\chi_m),


\end{equation}


т.\,к. $\chi_i^k=\chi_i$ для любых целых $i$ и $k$,


$M\chi_{i_1}\chi_{i_2}\dotsb \chi_{i_j}=M\chi_{1}\chi_{2}\dotsb \chi_{j}$


в силу симметричной зависимости $\chi_i.$


Необходимо вычислить моменты вида


$M(\chi_1\dotsb \chi_k).$





Нетрудно видеть, что $M(\chi_1\dotsc \chi_k)=P(\chi_1=1,\dotsc ,\chi_k=1).$


Зафиксируем первые $k$ пар $(t_1,u_1),\dotsc ,(t_k,u_k).$ Пусть все эти пары


истинны при натуральном совмещении, $(x_1,\dotsc ,x_k)$


упорядоченны по возрастанию $(t_1,\dotsc ,t_k)$, а $(y_1,\dotsc ,y_k)$


упорядоченны \tts{! было упорядоченныЕ}


$(u_1,\dotsc ,u_k)$. Тогда множества пар


$\{(x_1,y_1),\dotsc,(x_k,y_k)\}$ и $\{(t_1,u_1),\dotsc,(t_k,u_k)\}$ совпадают.


Посмотрим, как могут распределиться остальные $n-k$ пар. Ясно, что, сложив


вероятности возможных исходов, будем иметь искомую вероятность.





Чтобы облегчить описание процедуры подсчета, введем несколько


понятий.





Пусть пары $(x_i,y_i)$, $i\in 1,\dotsc ,k$, истинны при натуральном совмещении.


Точки с координатами $(x_i,y_i)$, $i\in1,\dotsc ,k$, назовем узлами. Каждый


узел прямыми $y=y_i$ и $x=x_i$ делит плоскость на четыре части.


Очевидно, для каждого узла число остальных пар, попавших в левую верхнюю


четверть, и в нижнюю правую должны совпадать. Это число для $i$-й


истинной пары будем называть $i$-м ящиком и обозначать $S_i$.


Далее, левую верхнюю четверть назовем для удобства $i$-м


верхним ящиком, правую нижнюю --- $i$-м нижним ящиком.





Каждой клетке сопоставим вероятность попадания в нее и


число пар, в нее попавших. Клетка --- прямоугольник,


ограниченный прямыми $y=y_i$, $y=y_{i+1}$, $x=x_i$, $x=x_{i+1}$, где


$x_0=y_0=-\infty$, $x_{k+1}=y_{k+1}=+\infty$.


Нумерацию строк будем вести снизу вверх, столбцов --- слева направо.





Для каждого ящика число пар, попавших в клетки, прилегающие к главной


диагонали (соединяющей левую нижнюю клетку с правой верхней), будем


вычислять как $S_i$ минус число пар, попавших в остальные клетки ящика.


Эти клетки назовем соответственно результирующими верхнего (нижнего)


ящика. Клетки ящика, не прилегающие к главной диагонали, назовем


элементами ящика. По аналогии с понятием ящика будем числу пар, попавших в


клетку, давать то же название, что и самой клетке.





Итак, пусть $k_{ij}$ --- элемент ящика. Символом $\sum_i^{(1)}$ будем


обозначать суммы элементов $i$-го верхнего ящика, символом $\sum_i^{(2)}$


--- $i$-го нижнего. Пусть, наконец, $p_{ij}$ --- вероятность попадания в


клетку $(i,j)$, $g=1-\sum\limits_{i\neq j}p_{ij}$, а $\sum^*$ --- число пар,


не попавших на главную диагональ.


Искомая вероятность равна


$$


\sum_{k_{ij},S_i}


\frac{


(n-k)!


\prod\limits_{| i-j | \ge 2}p_{ij}^{k_{ij}}


\prod\limits_i \Big(p_{i,i+1}^{S_i-\sum_i^{(2)}}\Big)


\Big(p_{i+1,i}^{S_i-\sum_i^{(1)}}\Big)


\Big(g^{n-k-\sum^*}\Big)}


{


\prod\limits_{| i-j| \ge 2}(k_{ij})! \,


\prod\limits_i \Big(S_i-\sum_i^{(1)}\Big)!\,


\Big(S_i-\sum_i^{(2)}\Big)!\,


\Big(n-k-\sum^*\Big)!},


$$


где суммирование проводится при неотрицательности всех чисел, стоящих в


знаменателе под знаком факториала; $k_{ij}$ исчезли при суммировании


(см.~[4]).





Легко видеть, что каждый член суммы равен


$$


\frac{1}{(2\pi i)^{k^2+k}} \int\limits_{\Gamma \subset C^{k^2+k}}


\frac{


\Big(g+\sum\limits_{i\neq j}p_{ij}u_{ij}\Big)^{n-k}du}


{


\prod\limits_{i\neq j}u_{ij}\prod\limits_{| i-j| \ge 2}


u_{ij}^{k_{ij}}


\prod\limits_i\Big(u_{i,i+1}^{S_i-\sum_i^{(1)}}


u_{i+1,i}^{S_i-\sum_i^{(2)}}\Big)


},


$$


где $\Gamma$ --- декартово произведение $k^2+k$ замкнутых контуров на 


комплексных плоскостях, содержащих точку $(0,0)$ внутри себя.


Если какая-то из степеней в знаменателе будет больше $n-k$,


то соответствующий интеграл обратится в нуль; поэтому суммирование по всем


переменным можно вести до бесконечности.


Теперь все суммы по $k_{ij}$ --- бесконечные геометрические прогрессии.


Выберем контур интегрирования так, чтобы эти прогрессии бесконечно


убывали.





Обозначив искомую вероятность при фиксированных


$(x_1,y_1),(x_2,y_2),\dotsc ,(x_k,y_k)$ через $P$, будем иметь


\begin{multline*}


P = \frac{1}{(2 \pi i)^{k^2+k}}\sum_{S_i}


\int\limits_{\Gamma \subset C^{k^2+k}}


\frac{


\Big( g+\sum\limits_{i \neq j} p_{ij} u_{ij}\Big)^{n-k}}


{


\prod\limits_i (u_{i,i+1}u_{i+1,i})^{S_i}


\prod\limits_{i\neq j} u_{ij} \prod\limits_{i-j>1}


\Big(1- \prod\limits^{i-1}_{m=j} u_{m+1,m}/u_{ij}\Big)} \times \\


%


\times \frac{


du}{\prod\limits_{i-j<-1}\Big(1-


\prod\limits^{j-1}_{m=i}u_{m,m+1}/u_{ij}\Big)}.


\end{multline*}


Проинтегрируем по $u_{ij}$, где $| i-j | >1$, затем просуммируем


по ящикам и проинтегрируем по $u_{i+1,i}$, т.\,е. по соответствующим


результирующим верхних ящиков. В итоге получим


$$


P=\frac{1}{(2 \pi i)^k}


\int\limits_{\Gamma \subset C^k}f(u) du,


$$


где


$$


f(u) = \frac{\bigg(g + \sum\limits_{i<j} \Big(p_{ij}\prod\limits_


{m=i}^{j-1} u_{m,m+1}+p_{ji}\prod\limits_{m=i}^{j-1}


u_{m,m+1}^{-1}\Big)\bigg)^{n-k}


}


{\prod\limits_{i=1}^k u_{i,j+1}}.


$$


Осталось осреднить $P$ по возможным значениям


$(x_1,y_1),(x_2,y_2),\dotsc , (x_k,y_k)$ и учесть,


что всего существует $k!$ перестановок


$(x_1,y_1), (x_2,y_2),\dotsc , (x_k,y_k)$.


Тогда для момента $M(\chi_1\dotsc\chi_k)$ имеем интегральное представление


\begin{equation}


\frac{k!}{(2\pi i)^k}


\idotsint\limits_{x_1<\dotsb <x_k}


\cdot


\idotsint\limits_{y_1<\dotsb <y_k}


\cdot


\int\limits_{\Gamma \subset C^k}


f(u) du \times 


\prod_{i=1}^k(f(x_i,y_i)dx_idy_i).


\end{equation}


Для нахождения $MN^k$ осталось применить (2.1).





\section{Асимптотика $MN^k$ при $n\to\infty$}





Как и в предыдущем пункте, займемся сначала моментами вида


$M(\chi_1\dotsb\chi_k)$, а затем воспользуемся формулой (2.1).





Для нахождения асимптотики $M(\chi_1\dotsb\chi_k)$ применим метод


перевала по переменным $u_{m,m+1}$ и метод Лапласа по переменным $y_i$.


Обозначим


$$


H=H(u,x,y)= g+\sum\limits_{i<j}


\bigg(p_{ij}\prod^{j-1}_{m=i} u_{m,m+1} +


p_{ji}\prod^{j-1}_{m=i} u^{-1}_{m,m+1}\bigg),


$$


где


\begin{gather*}


g=1-\sum_{i<j}(p_{ij}+p_{ji}),\\


u=(u_{12},u_{23},\dotsc,u_{k,k+1}), \quad x=(x_1,\dotsc,x_k), \quad


y=(y_1,\dotsc,y_k),


\end{gather*}


причем $y_1<y_2<\dotsb<y_k$, $x_1<x_2<\dotsb<x_k$.





Для нахождения асимптотики интеграла по замкнутому контуру $\Gamma$ необходимо


найти точку, в которой достигается минимакс


\begin{equation}


\min_{\Gamma \in G} \, \max_{u \in \Gamma} \re H(u,x,y),


\end{equation}


где $G$ --- множество всех контуров, сохраняющих значение интеграла.


В качестве контуров $G$ будем рассматривать окружности $u_{m,m+1}


=r_m\exp({i\psi_m})$, $m=1,\dotsc,k$. Тогда для нахождения асимптотики


интеграла (2.2) нужно найти точку $(y^*,u^*)$, в которой достигается минимакс


\begin{equation}


L(x)= \max_{y}\, \min_{r_m} \, \max_{\psi_m}


\re\, H(u,x,y).


\end{equation}





\begin{thm}


Минимакс в $(3.2)$ $L(x) \equiv 1$ $\forall x$ и достигается при 


$$


u_{m,m+1}^*=1, \quad m=1,\dotsc ,k; \quad


y_i^*(x_i)=F_2^{-1}(F_1(x_i)), \quad i=1,\dotsc ,k,


$$


где


$$


F_1(x)=\int_{-\infty}^x


\int_{-\infty}^{+\infty}f(x,y)dx\,dy,


\qquad


F_2(y)=\int_{-\infty}^{+\infty}


\int_{-\infty}^yf(x,y)dx\,dy.


$$


\end{thm}





\begin{pf}


Найдем сначала $\max\limits_{\psi_m} \re H(u,x,y)$. Имеем


\begin{multline*}


\re H(u,x,y)=\re


\bigg(g+\sum_{i<j}(p_{ij}r_i\dotsc r_{j-1}


\exp{(i(\psi_1+\dotsb +\psi_{j-1})))}+\\


%


+p_{ji}r_i^{-1}\dotsc r_{j-1}^{-1}


\exp(-i(\psi_1+\dotsb +\psi_{j-1}))\bigg) =\\


%


=g+\sum_{i<j}(p_{ij}r_i\dotsc r_{j-1}\cos(\psi_1+\dotsb +\psi_{j-1})+


p_{ji}r_i^{-1}\dotsc r_{j-1}^{-1}\cos(\psi_1+\dotsb +\psi_{j-1})),


\end{multline*}


и максимум достигается при $\psi_1=\psi_2=\dotsb=\psi_k=0.$





В точке минимакса ${{\partial H}\over {\partial y_1}}=0\;$, поэтому


\begin{multline*}


\frac{\partial H}{\partial y_1}=-(1-r_1)


\int_{-\infty}^{x_1} f(x,y_1)dx+


\bigg(1-\frac{1}{r_1}\bigg)\int_{x_1}^{x_2}f(x,y_1)dx+\\


%


+\bigg(1-\frac{1}{r_1r_2}\bigg)\int_{x_2}^{x_3}f(x,y_1)dx+\dotsb +


\bigg(1-\frac{1}{r_1\dotsc r_{k-1}}\bigg)\int_{x_k}^{+\infty}f(x,y_1)dx-\\


%


-\bigg(1-\frac{1}{r_2}\bigg)\int_{x_2}^{x_3}f(x,y_1)dx-


\bigg(1-\frac{1}{r_2r_3}\bigg)\int_{x_3}^{x_4}f(x,y_1)dx-\dotsb -


\bigg(1-\frac{1}{r_2\dotsc r_{k-1}}\bigg)\int_{x_k}^{+\infty}f(x,y_1)dx=\\


%


=(r_1-1)\bigg(\int_{-\infty}^{x_1}f(x,y_1)dx+


\frac{1}{r_1}\int_{x_1}^{x_2}f(x,y_1)dx+


\frac{1}{r_1r_2}\int_{x_2}^{x_3}f(x,y_1)dx+\dotsb +\\


%


+\frac{1}{r_1\dotsc r_{k-1}}\int_{x_k}^{+\infty}f(x,y_1)dx\bigg)=0.


\end{multline*}


Отсюда $r_1=1$.


Аналогично, из условия


$\frac{\partial H}{\partial y_i}=0$


следует $r_i=1$, $i=2,\dotsc ,k$.





Итак, минимакс достигается при $u_{m,m+1}^*=1$, $m=1,\dotsc ,k$, поэтому


$\frac{\partial H}{\partial u_{m,m+1}} \big|_{u^*=1}=0$.


Значит,


$$


\frac{\partial H}{\partial u_{m,m+1}} \bigg|_{u^*=1}=


\sum_{i<j}(p_{ij}-p_{ji})=\int_{-\infty}^{x_m}\int_{y_m}^{+\infty}


f(x,y)dx\,dy-\int_{x_m}^{+\infty}


\int_{-\infty}^{y_m} f(x,y)dx\,dy = 0.


$$


Поскольку $F_1(x)$ и $F_2(y)$ --- возрастающие непрерывные функции, то


$\forall x_m\in R$ $\exists\,z\in(0,1):F_1(x_m)=z$, а для


$z$\; $\exists\,y_m:F_2(y_m)=z$. Следовательно,


$F_1(x_m)=F_2(y_m)$, а $L(x)=1$.


\end{pf}





Теперь можно получить асимптотику интеграла (2.2) (см.~[6],


с.\,418, предложение 4.1, с.\,122, теорема~4.1) в виде


\begin{equation}


k! \idotsint\limits_{x_1<\dotsb <x_k}


\frac{


\prod\limits_i f(x_i,y_i^\ast)dx_i}


{


\sqrt{\det H(z,x,y)_{zz}''\big|\begin{Sb} z=1 \\ y=y^\ast\end{Sb}


\det H(r^\ast,x,y)_{xy}''\big|\begin{Sb} r^\ast=1 \\ y=y^\ast \end{Sb}}


},


\end{equation}


где $r^\ast(x,y)$ --- радиус перевального контура, а $y^\ast:


F_1(x_m)=F_2(y_m^\ast)$ $\forall m$.


Заметим, что в точке $(r^\ast=1,\ y=y^\ast)$ $p_{ij}=p_{ji}$ для любых $i$ и


$j$.





\section{Выражение для $DN$}





Получим выражение


$DN=MN^2-(MN)^2=nM\chi_1+n(n-1)M\chi_1\chi_2-n^2(M\chi_1)^2$.


Найдем $\lim\limits_{n \to +\infty}n(n-1)M\chi_1\chi_2$.


В соответствии с (3.3) имеем


$$


n^2M\chi_1\chi_2 \sim 2! \iint\limits_{x_1<x_2}


\frac{


f(x_1,y_1^\ast(x_1)) f(x_2,y_2^\ast(x_2)) dx_1dx_2


}


{


\sqrt{\det H(z,x,y)_{zz}''\big|\begin{Sb} z=1 \\ y=y^\ast\end{Sb}


\det H(r^\ast,x,y)_{xy}''\big|\begin{Sb} r^\ast=1 \\ y=y^\ast\end{Sb}}


}.


$$





\begin{thm} Имеет место равенство


\begin{equation}


\sqrt{\det H(z,x,y)_{zz}''\big|\begin{Sb} z=1 \\ y=y^\ast\end{Sb}


\det H(r^\ast,x,y)_{xy}''\big|\begin{Sb} r^\ast=1 \\ y=y^\ast\end{Sb}}


=f_2(y^\ast_1)f_2(y_2^\ast)>0.


\end{equation}


\end{thm}





\begin{pf}


Для $k=2$ имеем


$$


H=g+p_{12}z_1+p_{21}\frac{1}{z_1}+p_{23}z_2+p_{32}


\frac{1}{z_2}+p_{13}z_1z_2+p_{31}\frac{1}{z_1z_2}


$$


и


$$


\det H(z,x,y)_{zz}''\big|\begin{Sb} z=1 \\ y=y^\ast \end{Sb}=


4(p_{12}p_{23}+p_{12}p_{13}+p_{23}p_{13}).


$$





Для получения второго определителя найдем производные вида


$(r_i^{\ast})'_{y_j}$ в точке $(r^\ast=1,\ y=y^\ast)$. Так как


$r^\ast$ --- радиус перевального контура, то


$H(r^\ast,x,y)'_{r_i^\ast}=0$, т.\,е.


$$


p_{12}-\frac{p_{21}}{(r_1^\ast)^2}+p_{13}r_2^\ast-


\frac{p_{31}}{(r_1^\ast)^2r^\ast_2}=0,\quad


p_{23}-\frac{p_{32}}{(r_2^\ast)^2}+p_{13}r_2^\ast-


\frac{p_{31}}{(r_2^\ast)^2r^\ast_1}=0.


$$





Продифференцируем каждое из этих уравнений по $y_1$


$$


f_2(y_1)+


2r_{1\, y_1}^{\ast'}(p_{12}+p_{13})+2p_{13}


r^{\ast'}_{2 y_1}=0; \quad


2r_{1 \, y_1}^{\ast'}(p_{23}+p_{13})+


2p_{13}r^{\ast'}_{1\,y_1}=0.


$$


Отсюда


\begin{equation}


r_{1 y_1}^{\ast'}=-\frac{f_2(y_1)}{2}\frac{p_{23}+p_{13}}


{p_{23}p_{13}+p_{12}p_{13}+p_{23}+p_{12}}; \quad


r_{2 y_1}^{\ast'}=\frac{f_2(y_1)}{2}\frac{p_{13}}


{p_{23}p_{13}+p_{12}p_{13}+p_{23}+p_{12}}.


\end{equation}


Аналогично,


\begin{equation}


r_{1 y_2}^{\ast'}=\frac{f_2(y_2)}{2}\frac{p_{13}}


{p_{23}p_{13}+p_{12}p_{13}+p_{23}p_{12}}, \quad


r_{2 y_2}^{\ast'}=-\frac{f_2(y_2)}{2} \frac{p_{13}+p_{12}}


{p_{23}p_{13}+p_{12}p_{13}+p_{23}p_{12}}.\tag{4.2$'$}


\end{equation}





Введем обозначения


$r_{i y_j}^{\ast'}=a_{ij}$; $(p_{12}-p_{21})_{y_j}'=b_{1j}$;


$(p_{13}-p_{31})'_{y_j}=b_{2j}$;


$(p_{23}-p_{32})'_{y_j}=b_{3j}$.


Учитывая, что $b_{21}=-b_{31}=f_2(y_1)-b_{11}$;


$b_{22}=-b_{12}=f_2(y_2)-b_{32}$, имеем


\begin{align}


H''_{y_1y_1}&=2a_{11}b_{11}+2a_{21}b_{31}+2a_{11}^2p_{12}+2


a^2_{21}p_{23}+


2b_{21}(a_{11}+a_{21})+2p_{13}(a_{11}+a_{21})^2=\notag\\


%


&\qquad=2a_{11}f_2(y_1)+2a_{11}^2p_{12}+2a_{21}^2p_{23}+


2p_{13}(a_{11}+a_{21})^2,\notag \\


%


H''_{y_2y_2}&=2a_{22}f_2(y_2)+2a_{12}^2p_{12}+2a_{22}^2p_{23}+


2p_{13}(a_{12}+a_{22})^2,\\


%


H''_{y_1y_2}&=a_{12}f_2(y_1)+a_{12}^2f_2(y_2)+


2a_{11}a_{12}p_{12}+2a_{21}a_{22}p_{23}+2p_{13}(a_{12}+


a_{22})(a_{11}+a_{21}).\notag


\end{align}





Подставив (4.2), (4.2$'$) в (4.3), получим


\begin{multline*}


H_{y_1y_1}''H_{y_2y_2}''-(H_{y_1y_2}'')^2=


f_2^2{(y_1)}f_2^2{(y_2)}\bigg(a_{11}a_{22}\frac{1}{f_2(y_1)f_2(y_2)}-


\bigg(\frac{a_{12}}{f_2(y_2)}+\frac{a_{21}}{f_2(y_1)}\bigg)^2\bigg)=\\


%


=f_2^2{(y_1)}f_2^2{(y_2)}


\frac{1}{4(p_{12}p_{13}+p_{12}p_{23} +p_{23}p_{13})},


\end{multline*}


отсюда вытекает (4.1).


\end{pf}





Теперь асимптотика $M\chi_1 \chi_2$ значительно упрощается. При


$n\to+\infty$ имеем


\begin{multline*}


n^2M\chi_1 \chi_2 \sim 2\iint\limits_{x_1<x_2}


\frac{f(x_1,y_1^*)f(x_2,y_2^*)}


{f_2(y_1^*)f_2(y_2^*)}dx_1dx_2=\\


%


=\int_{-\infty}^{+\infty} \int_{-\infty}^{+\infty}


\frac{f(x_1,y_1^*)f(x_2,y_2^*)}{f_2(y_1^*)f_2(y_2^*)}dx_1dx_2=


\bigg(\int_{-\infty}^{+\infty}


\frac{f(x_1,y_1^*)}{f_2(y_1^*)}


dx_1 \bigg)^2 = (nM\chi_1)^2.


\end{multline*}


В итоге


$$


DN \sim MN \sim \int_{-\infty}^{+\infty}\frac{f(x_1,y_1^*)}


{f_2(y_1^* )}dx_1=\lambda>0.


$$





\section{Вид распределения числа совпадений}





{\bf Предложение} ([7], с.\,265). {\it


Пусть случайные величины $\chi_1,\dotsc,\chi_n$ принимают только


значения $0$ и $1$ и симметрично зависимы.


Тогда


\begin{itemize}


\item[1.] моменты $C_{nk}=M(\chi_1\dotsb\chi_k)$ образуют вполне


монотонную последовательность $($т.\,е. знаки конечных разностей элементов


последовательности чередуются) для каждого $n;$


\item[2.] существует случайная величина $\lambda_n$, распределение


которой сосредоточено на решетке $j/n$, $j=0,\dotsc,n$, такая, что


$M\lambda_n^k=C_{nk};$


\item[3.] случайные величины $N_n$ распределены по рандомизированному


биномиальному закону


$$


P(N_n=k)=\sum_{l=0}^{n}C_n^k\bigg(\frac{l}{k}\bigg)^k


\bigg(1-\frac{l}{k}\bigg)^{n_k}


P\bigg(\lambda_n=\frac{l}{k}\bigg).


$$


\end{itemize}


}





\begin{thm}


Пусть $N_n$ --- число совпадений для объема выборки $n$,


$$


C_{nk}=M(\chi_1\dotsb \chi_k), \quad S_k=\lim_{n \to +\infty}n^kC_{nk}.


$$


Тогда


\begin{itemize}


\item[а)] существует случайная величина $\lambda$, распределение которой


сосредоточено на решетке $j/n$, $j=0,\dotsc,n$, такая, что


$M\lambda_n^k=C_{nk}$, и


$p(N_n=k)=\sum\limits_{l=0}^n C_n^k l^k(1-l)^{n-k}P


\big(\lambda_n=\frac{l}{k}\bigr);$


\item[б)] для любого частичного слабого предела $N$ последовательности


$N_n$ существует


$\mu \overset{\mathrm D}{=}{} \lim\limits_{n \to +\infty}


n\lambda_n : M\mu^k=S_k$


и распределение $N$ --- рандомизированное по параметру распределение Пуассона


$$


P(N=k)=\sum_{l=0}^{+\infty}\frac{e^{-l}l^k}{k!}P(\mu=l).


$$


\end{itemize}


\end{thm}





\begin{pf}


Легко видеть, что утверждение а) прямо следует из пп.~2, 3


предложения, т.\,к. $\chi_1,\chi_2,\dotsc ,\chi_k$ симметрично зависимы.





б) Пусть $\mu_n=n\lambda_n$, $F_n$ --- функция распределения $\mu_n$.


Легко показать, что последовательность $\mu_n$ стохастически ограничена.


Действительно, из неравенства Чебышева


$$


P(\mu_n \ge t) \leq \frac{1}{t^2}M(\mu_n^2).


$$


Для фиксированного $\varepsilon^*(n>n^*)


\rightarrow P(\mu_n \ge t) \leq \frac{1}{t^2}( S_2+\varepsilon^*).$


Тогда $\forall (\varepsilon>0,\ n>n^*)$ существует $a=


\sqrt{S_2+\varepsilon^*/\,\varepsilon}\, : P(\mu_n \ge a) \leq


\varepsilon$, что и означает стохастическую ограниченность


последовательности $\mu_n$.


Теперь из теоремы Хелли о выборе ([7], с.\,307) ясно, что если


последовательность $F_n$ и не имеет предела, то все ее частичные пределы


будут функциями распределения и будут обладать тем свойством, что их


$k$-й момент равен $S_k$.


Выберем любой из этих пределов. Пусть $\mu$ -- случайная величина, ему


соответствующая. Далее утверждение~б) становится очевидным.


\end{pf}





\section{Асимптотическое распределение числа совпадений}





Для нахождения предельного распределения числа совпадений воспользуемся


методом моментов. Прежде всего методом математической индукции выведем


рекуррентные соотношения для коэффициентов $c_{jk}$ в (2.1).





При $k=1$, очевидно, $c_{11}=1$.


Предположим, что известны $c_{jk}$, $j=1,\dotsc ,k$, и


\begin{multline*}


MN_n^{k+1}=M\bigg(\sum_{j=1}^k c_{jk}


\underset{i_1\neq i_2\neq \dotsb \neq i_j}


{\sum\dotsb \sum}{} \chi_{i_1}\dotsc \chi_{i_j}\bigg)


\bigg(\sum_{i=1}^n\chi_i\bigg)= \\


%


=M\bigg(\sum_{j=1}^k c_{jk}


\underset{i_1\neq \dotsb \neq i_{j+1}}


{\sum\dotsb \sum}{}\chi_{i_1}\dotsc \chi_{i_{j+1}}+


\sum_{j=1}^k j c_{jk}


\underset{i_1\neq \dotsb \neq i_j}


{\sum\dotsb \sum}\chi_{i_1}\dotsc \chi_{i_j}\bigg).


\end{multline*}


Таким образом, $c_{1,k+1}=1$, $c_{j,k+1}=c_{j-1,k}+j c_{jk}$.





Предположим, что $с_{jk}$ суть коэффициенты полинома


$\Phi_k(\lambda)=\sum\limits_{j=1}^k c_{jk} \lambda^k$, тогда


\begin{equation}


\begin{array}{c}


\Phi_1(\lambda)=1,\\[2mm]


\Phi_{k+1}(\lambda)=\lambda(\Phi_k(\lambda)+\Phi_k'(\lambda)), 


\quad k =2,3, \dotsc 


\end{array}


\end{equation}


Из (2.1) и (3.3) имеем


$\lim\limits_{n\to\infty} MN_n^k=\sum\limits_{j=1}^k c_{jk} \lambda^j=


\Phi_k(\lambda)$, 


где $\lambda=\lim\limits_{n\to\infty} n M \chi_1$.





Начальные моменты случайной величины $W$, подчиненной распределению Пуассона,


удовлетворяют (6.1). Действительно, $MW=\lambda=\Phi_1(\lambda)$. 


Предположим $MW^k=\Phi_k(\lambda)$.


Это означает, что


$\phi^{(k)}(t)=(M e^{itw})^{(k)}= e^{\lambda(e^{it}-1)}


(\Phi_k(\lambda e^{it}))i^k=e^{\lambda(e^{it}-1)}


\sum\limits_{j=1}^k c_{jk}\lambda^j e^{ijt} i^k$.


Дифференцируя последнее равенство, имеем


$\phi^{(k+1)}(t)=e^{\lambda(e^{it}-1)}


\Big(\sum\limits_{j=1}^k \lambda e^{it} i c_{jk}


\lambda^j e^{ijt}+i j c_{jk} \lambda^j e^{ijt}\Big)i^k=


e^{\lambda(e^{it}-1)} \Phi_{k+1}(\lambda e^{it}) i^{k+1}$.


Следовательно, $MW^{k+1}=\Phi_{k+1}(\lambda)$.





Так как распределение Пуассона однозначно определяется своими моментами,


предельное распределение $N_n$ пуассоновское.





Заметим, что при $\lambda<\infty$ этот результат следует из теоремы 3,


поскольку


\begin{gather*}


M\mu=\lim_{n\to\infty} M\mu_n = \lambda<+\infty,\\


M\mu_n^2={(M\mu_n)}^2+{\mathrm o}({(M\mu_n)}^2)={(M\mu_n)}^0+


{\mathrm o}(1),


\end{gather*}


т.\,к. $n^2M \chi_1 \chi_2 \sim {(n M \chi_1)}^2$ и 


$D\mu=\lim\limits_{n\to\infty} D\mu_n = 0.$





Покажем, что при $\lambda\to\infty$ распределение числа совпадений $N$


асимптотически нормально. Пусть $\{ W_n\}$ --- последовательность 


случайных величин, подчиненных закону Пуассона


с параметрами $\lambda_n=nM\chi_1$.


Случайные величины $N_n$ и $W_n$ определены на одном и том же вероятностном


пространстве $\bold Z$ и сходятся по распределению к одной и той же


случайной величине $W$, поэтому


$N_n-W_n \mathop{\rightarrow}\limits^p_{n \to 0}0.$


Имеем $N_n=W_n+(N_n-W_n)$, поскольку при $\lambda\to\infty$


распределение $W_n$ асимптотически нормально, таким же 


является и предельное распределение $N_n$, т.е.


$$


P\bigg(\frac{N_n-MN_n}{\sqrt{DN_n}}<x \bigg)\to \Phi(x).


$$





\begin{thebibliography}{9}





\bibitem{1}


De Groot M.H., Feder P.I., Goel P.K. {\em Matchmaking}


// Ann. Math. Stat. -- 1971. --  V.\,42. -- \No.\,2. -- P.\,578--593.


\bibitem{2}


Chew M.C. {\em On pairing observations from a distribution with monotone


likehood ratio} // Ann. Math. Stat. --  1973. -- V.\,1.


-- \No3. -- P.\,433--445.


\bibitem{3}


De Groot M.H., Goel P.K. {\em The matching problem for multivariate


normal data} // Sankhya. -- 1976. -- V.\,38. -- \No\,1. -- P.\,14--29.


\bibitem{4}


Золотухина Л.А., Латышев К.П. {\em Асимптотическое представление


среднего числа совпадений элементов вариационных рядов двумерной


выборки} // Зап. научн. семин. ЛОМИ. -- 1978. -- T.\,79. -- C.\,4--10.


\bibitem{5}


Латышев К.П., Золотухина Л.А. {\em Интегральное представление дисперсии


числа правильно угаданных пар в задаче о совпадении двумерной выборки}


// Зап. научн. семин. ЛОМИ. -- 1982. -- T.\,136. -- C.\,113--120.


\bibitem{6}


Федорюк М.В. {\em Асимптотика. Интегралы и ряды}.


-- М.: Наука, 1987. -- 361~с.


\bibitem{7}


Феллер В.И. {\em Введение в теорию вероятностей и ее применения}. T.\,2.


-- М.: Мир, 1984. -- 751~с.





\end{thebibliography}





\vskip 0.5cm





\postup{\it Санкт-Петербургский государственный}{$\qquad$\it Поступили}


\postup{\it морской технический университет}{\it первый вариант $11.04.1997$}


\postup{}{\it окончательный вариант $29.10.1998$}





\label{end}


\end{document}





